\documentclass[10pt,a4paper]{article}
\usepackage[margin=2.1cm]{geometry}
\usepackage{fontspec}
\usepackage{xeCJK}
\setCJKmainfont{Noto Serif CJK KR}
\setCJKsansfont{Noto Sans CJK KR}
\setCJKmonofont{Noto Sans Mono CJK KR}
\xeCJKsetup{CJKspace=true}
\usepackage{amsmath,amssymb,mathtools}
\usepackage{booktabs,tabularx,multirow,array,longtable}
\usepackage{enumitem}
\usepackage{xcolor}
\usepackage{caption}
\usepackage{tikz}
\usetikzlibrary{positioning,arrows.meta,fit,backgrounds,calc,shapes.geometric,decorations.pathreplacing}
\usepackage[colorlinks=true]{hyperref}
\linespread{1.15}
\setlength{\parskip}{3pt}
\setlist{itemsep=1pt,topsep=2pt}

% ---------- notation macros ----------
\newcommand{\G}{\mathcal{G}}
\newcommand{\KS}{\mathcal{K}_S}
\newcommand{\Ont}{\mathcal{O}}
\newcommand{\Ctx}{\mathcal{C}}
\newcommand{\Tcmd}{\mathcal{T}_{\mathrm{cmd}}}
\newcommand{\Treq}{\mathcal{T}_{\mathrm{req}}}
\newcommand{\Tst}{\mathcal{T}_{\mathrm{state}}}
\newcommand{\Tadv}{\mathcal{T}_{\mathrm{adv}}}
\newcommand{\Yset}{\mathcal{Y}}
\newcommand{\Vdom}{\mathcal{V}_{\mathrm{dom}}}
\newcommand{\KR}{\mathcal{K}_R}
\newcommand{\Cfg}{\Theta_{\mathrm{kg}}}
\newcommand{\Ivl}{\mathbb{I}}
\newcommand{\ATCO}{\textsf{ATCO}}
\newcommand{\PILOT}{\textsf{PILOT}}
\newcommand{\cmd}[1]{\textsf{\small #1}}
\newcommand{\ch}[1]{\textbf{\textcolor{red!60!black}{#1}}}
\newcommand{\md}[1]{\textbf{\textcolor{blue!60!black}{#1}}}

\definecolor{perc}{RGB}{225,238,250}
\definecolor{comp}{RGB}{232,245,228}
\definecolor{proj}{RGB}{252,240,222}
\definecolor{res}{RGB}{240,236,248}
\pagestyle{empty}
\begin{document}
\section*{KG 스키마 그림 후보 비교}
네 후보 모두 \textbf{같은 KG}를 그린 것이다(노드 종류 Transmission·Instruction은 각각 하나). 다른 점은 그리는 방식과, D만 복창 간선을 지시 단위로 바꾸는 구조 변경이라는 점이다.
\begin{center}\small
\begin{tabularx}{\textwidth}{@{}l>{\raggedright\arraybackslash}X>{\raggedright\arraybackslash}X>{\raggedright\arraybackslash}X@{}}
\toprule
후보 & 그리는 방식 & 장점 & 단점\\\midrule
A (현재) & 예시식: Transmission·Instruction을 여러 번 그림 & 관계가 구체적으로 보임 & 노드 종류가 나뉜 것처럼 보임\\
B & 스키마식: 종류마다 한 번, 같은 종류끼리 관계는 자기 화살표(self-loop) & 노드 종류와 관계가 한눈에 보임 & 시간 흐름(τ, 순서)은 안 보임\\
C & B + KATL 실제 예시 시간축 & 구조(B)와 시간·복창(예시)을 함께 보여줌 & 그림이 2개\\
D & B와 같되 READBACK\_OF를 \emph{지시끼리} 연결 (구조 변경) & 어느 지시가 복창됐는지 그래프에서 바로 보임 & 코드 수정 필요, 간선 수 증가\\
\bottomrule
\end{tabularx}
\end{center}
\newpage
\section*{후보 A: 현재 그림 (예시식)}
관제사 송신과 조종사 송신, 지시들을 \textbf{각각 따로} 그려 관계를 보여준다. 실제 KG에서는 Transmission·Instruction이 각각 한 종류이고, 관제사/조종사는 속성 \texttt{speaker}로 구분된다.
\begin{center}
\resizebox{0.95\textwidth}{!}{%
\begin{tikzpicture}[
  font=\small,
  dyn/.style={draw, thick, rounded corners=2pt, fill=comp, minimum height=0.8cm, minimum width=2.4cm, align=center},
  sta/.style={draw, thick, fill=res, minimum height=0.8cm, minimum width=2.2cm, align=center},
  e/.style={-{Stealth[length=2.2mm]}, thick},
  el/.style={font=\scriptsize\sffamily, fill=white, inner sep=1pt}
]
\node[dyn] (tx) at (0,0) {\textbf{Transmission}\\{\scriptsize (ATCO) $s, w$}};
\node[dyn] (tx2) at (0,-3.2) {\textbf{Transmission}\\{\scriptsize (PILOT)}};
\node[sta] (role) at (-4.2,0) {\textbf{Role}\\{\scriptsize ATCO / PILOT}};
\node[dyn] (fl) at (4.6,0) {\textbf{Flight}\\{\scriptsize $c$ (예: DAL1489)}};
\node[sta] (al) at (8.6,0) {\textbf{Airline}\\{\scriptsize icao, country (예: DAL, US)}};
\node[dyn] (ins) at (4.6,-3.2) {\textbf{Instruction}\\{\scriptsize $\kappa,v,q,\tau$}};
\node[sta] (rw) at (8.6,-3.2) {\textbf{Runway / Taxiway}\\{\scriptsize / Facility}};
\node[sta] (ap) at (8.6,-5.6) {\textbf{Airport}};
\node[dyn] (cm) at (4.6,2.2) {\textbf{CallsignMention}\\{\scriptsize $\hat c$ (원 표면형)}};
\node[dyn] (ins2) at (4.6,-5.8) {\textbf{Instruction}\\{\scriptsize 다음 관제사 지시 (같은 항공편)}};
\node[dyn] (pins) at (0,-5.8) {\textbf{Instruction}\\{\scriptsize 조종사가 복창한 지시}};

\draw[e] (tx) -- node[el,above]{SPOKEN\_BY} (role);
\draw[e] (tx) -- node[el,above]{ADDRESSES} (fl);
\draw[e] (fl) -- node[el,above]{OPERATED\_BY} (al);
\draw[e] (tx) -- node[el,sloped,above]{CONTAINS} (ins);
\draw[e] (ins) -- node[el,above]{TARGET} (rw);
\draw[e] (rw) -- node[el,right]{PART\_OF} (ap);
\draw[e] (fl) -- node[el,right]{ALIAS} (cm);
\draw[e, red!60!black] (tx2) -- node[el,left,text=red!60!black]{READBACK\_OF \{status\}} (tx);
\draw[e] (ins) -- node[el,right]{NEXT \{$\Delta t$\}} (ins2);
\draw[e] (tx2) -- node[el,right]{CONTAINS} (pins);

\node[font=\scriptsize, align=left, anchor=north west] at (-5.6,-7.0) {
\tikz\node[fill=comp,draw,minimum width=4mm,minimum height=3mm]{}; 동적(사건) 계층: 유효 구간 $\tau$, 근거 송신 prov 보유\\[2pt]
모든 간선: 적재 시각 $t'$ (\texttt{ingested\_at})\\[2pt]
\tikz\node[fill=res,draw,minimum width=4mm,minimum height=3mm]{}; 정적 계층 $\KS$ (atmonto IRI 매핑)};
\end{tikzpicture}}
\end{center}
\newpage
\section*{후보 B: 스키마식 (종류마다 한 번 + 자기 화살표)}
Transmission·Instruction을 \textbf{한 번씩만} 그리고, 같은 종류끼리의 관계는 자기 자신으로 돌아오는 화살표로 그린다. 관제사/조종사는 상자 안 \texttt{speaker} 속성으로 표시한다.
\begin{itemize}\itemsep0pt
\item READBACK\_OF: 조종사 송신 $\to$ 관제사 송신 (복창 관계)
\item NEXT: 같은 항공편의 관제사 지시 $\to$ 다음 관제사 지시 (시간 순서)
\end{itemize}
\begin{center}
\resizebox{0.95\textwidth}{!}{%
\begin{tikzpicture}[
  font=\small,
  dyn/.style={draw, thick, rounded corners=2pt, fill=comp, minimum height=0.8cm, minimum width=2.4cm, align=center},
  sta/.style={draw, thick, fill=res, minimum height=0.8cm, minimum width=2.2cm, align=center},
  e/.style={-{Stealth[length=2.2mm]}, thick},
  el/.style={font=\scriptsize\sffamily, fill=white, inner sep=1pt}
]
\node[dyn] (tx) at (0,0) {\textbf{Transmission}\\{\scriptsize $s, w$, speaker}};
\node[sta] (role) at (-4.4,0) {\textbf{Role}\\{\scriptsize ATCO / PILOT}};
\node[dyn] (fl) at (4.8,0) {\textbf{Flight}\\{\scriptsize $c$ (예: DAL1489)}};
\node[sta] (al) at (9.0,0) {\textbf{Airline}\\{\scriptsize icao, country}};
\node[dyn] (cm) at (4.8,2.2) {\textbf{CallsignMention}\\{\scriptsize $\hat c$ (원 표면형)}};
\node[dyn] (ins) at (0,-3.0) {\textbf{Instruction}\\{\scriptsize $\kappa,v,q,\tau$, speaker}};
\node[sta] (rw) at (4.8,-3.0) {\textbf{Runway / Taxiway}\\{\scriptsize / Ramp / Facility}};
\node[sta] (ap) at (9.0,-3.0) {\textbf{Airport}};
\draw[e] (tx) -- node[el,above]{SPOKEN\_BY} (role);
\draw[e] (tx) -- node[el,above]{ADDRESSES} (fl);
\draw[e] (fl) -- node[el,above]{OPERATED\_BY} (al);
\draw[e] (fl) -- node[el,right]{ALIAS} (cm);
\draw[e] (tx) -- node[el,right]{CONTAINS} (ins);
\draw[e] ([yshift=2mm]ins.east) -- node[el,above]{TARGET} ([yshift=2mm]rw.west);
\draw[e] ([yshift=-2mm]ins.east) -- node[el,below]{VIA} ([yshift=-2mm]rw.west);
\draw[e] (rw) -- node[el,above]{PART\_OF} (ap);
\draw[e, red!60!black] (tx.north west) to[out=110,in=180,looseness=1.2] ($(tx.north)+(-0.2,1.0)$) node[el,above,text=red!60!black,align=center]{READBACK\_OF \{status\}\\{\scriptsize PILOT 송신 $\to$ 직전 ATCO 송신: 같은 항공편·주파수, $0<\Delta t\le 20$ s}} to[out=0,in=70,looseness=1.2] (tx.north);
\draw[e, teal!60!black] (ins.west) to[out=200,in=270,looseness=1.3] ($(ins.west)+(-1.6,0.6)$) node[el,left,text=teal!60!black,align=right]{NEXT \{$\Delta t$\}\\{\scriptsize 같은 항공편의}\\{\scriptsize 다음 관제사 지시}\\{\scriptsize (시간순, $\Delta t\ge 0$)}} to[out=90,in=160,looseness=1.3] ([yshift=3mm]ins.west);
\end{tikzpicture}}
\end{center}
\newpage
\section*{후보 C: 스키마(B) + 실제 예시 시간축 (추천)}
위는 후보 B의 스키마(\emph{어떤 노드와 관계가 있나}), 아래는 KATL 지상관제 DAL1489의 실제 송신으로 그린 시간축(\emph{실제로 어떻게 연결되나})이다. 아래 그림이 기존 SFO 시간 모델 그림을 대신한다.
\begin{center}
\resizebox{0.8\textwidth}{!}{%
\begin{tikzpicture}[
  font=\small,
  dyn/.style={draw, thick, rounded corners=2pt, fill=comp, minimum height=0.8cm, minimum width=2.4cm, align=center},
  sta/.style={draw, thick, fill=res, minimum height=0.8cm, minimum width=2.2cm, align=center},
  e/.style={-{Stealth[length=2.2mm]}, thick},
  el/.style={font=\scriptsize\sffamily, fill=white, inner sep=1pt}
]
\node[dyn] (tx) at (0,0) {\textbf{Transmission}\\{\scriptsize $s, w$, speaker}};
\node[sta] (role) at (-4.4,0) {\textbf{Role}\\{\scriptsize ATCO / PILOT}};
\node[dyn] (fl) at (4.8,0) {\textbf{Flight}\\{\scriptsize $c$ (예: DAL1489)}};
\node[sta] (al) at (9.0,0) {\textbf{Airline}\\{\scriptsize icao, country}};
\node[dyn] (cm) at (4.8,2.2) {\textbf{CallsignMention}\\{\scriptsize $\hat c$ (원 표면형)}};
\node[dyn] (ins) at (0,-3.0) {\textbf{Instruction}\\{\scriptsize $\kappa,v,q,\tau$, speaker}};
\node[sta] (rw) at (4.8,-3.0) {\textbf{Runway / Taxiway}\\{\scriptsize / Ramp / Facility}};
\node[sta] (ap) at (9.0,-3.0) {\textbf{Airport}};
\draw[e] (tx) -- node[el,above]{SPOKEN\_BY} (role);
\draw[e] (tx) -- node[el,above]{ADDRESSES} (fl);
\draw[e] (fl) -- node[el,above]{OPERATED\_BY} (al);
\draw[e] (fl) -- node[el,right]{ALIAS} (cm);
\draw[e] (tx) -- node[el,right]{CONTAINS} (ins);
\draw[e] ([yshift=2mm]ins.east) -- node[el,above]{TARGET} ([yshift=2mm]rw.west);
\draw[e] ([yshift=-2mm]ins.east) -- node[el,below]{VIA} ([yshift=-2mm]rw.west);
\draw[e] (rw) -- node[el,above]{PART\_OF} (ap);
\draw[e, red!60!black] (tx.north west) to[out=110,in=180,looseness=1.2] ($(tx.north)+(-0.2,1.0)$) node[el,above,text=red!60!black,align=center]{READBACK\_OF \{status\}\\{\scriptsize PILOT 송신 $\to$ 직전 ATCO 송신: 같은 항공편·주파수, $0<\Delta t\le 20$ s}} to[out=0,in=70,looseness=1.2] (tx.north);
\draw[e, teal!60!black] (ins.west) to[out=200,in=270,looseness=1.3] ($(ins.west)+(-1.6,0.6)$) node[el,left,text=teal!60!black,align=right]{NEXT \{$\Delta t$\}\\{\scriptsize 같은 항공편의}\\{\scriptsize 다음 관제사 지시}\\{\scriptsize (시간순, $\Delta t\ge 0$)}} to[out=90,in=160,looseness=1.3] ([yshift=3mm]ins.west);
\end{tikzpicture}}

\vspace{6mm}
\resizebox{\textwidth}{!}{%
\begin{tikzpicture}[font=\small, x=0.13cm, y=1cm]
\draw[-{Stealth}, thick] (18,0) -- (102,0) node[right]{$t$ [s]};
\foreach \x in {20,30,...,100} { \draw (\x,0.08) -- (\x,-0.08) node[below,font=\scriptsize]{\x}; }
% transmissions (short labels; contents listed below the figure)
\fill[blue!25] (24.1,0.85) rectangle (30.4,1.15); \node[anchor=south,font=\scriptsize] at (27.2,1.15) {\#0005 ATCO};
\fill[orange!40] (31.7,1.85) rectangle (37.2,2.15); \node[anchor=west,font=\scriptsize] at (37.6,2.0) {\#0006 PILOT (복창)};
\fill[blue!25] (88.5,0.85) rectangle (91.2,1.15); \node[anchor=south,font=\scriptsize] at (89.8,1.15) {\#0022 ATCO};
% readback
\draw[red!60!black, thick, -{Stealth}] (33,1.83) to[out=-110,in=30] (29.5,1.5);
\node[font=\scriptsize, red!60!black, anchor=west] at (38,1.55) {READBACK\_OF \{MISMATCH\}};
\draw[red!60!black, thin] (37.9,1.55) -- (33.2,1.6);
% next
\draw[teal!60!black, thick, dashed, -{Stealth}] (30.4,0.75) to[out=-10,in=190] (88.5,0.75);
\node[font=\scriptsize, teal!60!black, fill=white, inner sep=1pt] at (72,0.3) {NEXT \{$\Delta t$ = 64.4 s\}};
% validity bars
\fill[green!45!black!60] (24.1,-0.95) rectangle (88.5,-0.7); \node[anchor=east,font=\scriptsize] at (23.5,-0.82) {$\tau$(TAXI\_VIA)};
\node[font=\scriptsize, anchor=west] at (88.8,-0.82) {MONITOR가 끝냄};
\fill[green!25] (24.1,-1.45) rectangle (30.4,-1.2); \node[anchor=east,font=\scriptsize] at (23.5,-1.33) {$\tau$(FOLLOW)};
\fill[green!45!black!60] (88.5,-1.95) rectangle (101,-1.7); \node[anchor=east,font=\scriptsize] at (88.0,-1.83) {$\tau$(MONITOR): 다음 지시가 없어 최대 1800 s까지};
% snapshot
\draw[dashed] (50,-2.2) -- (50,2.5) node[above,font=\scriptsize]{$\G|_{50}$};
\node[font=\scriptsize, anchor=north west] at (50.6,-2.15) {$\G|_{50}$: DAL1489는 TAXI\_VIA(M 경유 M2) 유효};
\end{tikzpicture}}
\end{center}
{\small KATL 전체(Whisper)에서 연결된 복창 404건 중 351건이 관제사 송신 끝난 뒤 5\,s 안, 95\%가 12.4\,s 안에 시작했다. 아래 예의 \#0006도 \#0005 끝난 뒤 1.3\,s 만에 시작한다.

\begin{tabularx}{\textwidth}{@{}ll>{\raggedright\arraybackslash}X@{}}
\#0005 & 24.1 s, ATCO & ``Delta fourteen eighty nine pass behind \dots\ nine left to Mike two taxi via Mike'' $\to$ FOLLOW, TAXI\_VIA (M 경유 M2), TRAFFIC\_INFO\\
\#0006 & 31.7 s, PILOT & ``Nine left Mike two via Mike \dots\ Delta fourteen eighty nine'' $\to$ TAXI\_VIA, FOLLOW (복창)\\
\#0022 & 88.5 s, ATCO & ``Delta fourteen eighty nine monitor twenty three eighty five'' $\to$ MONITOR\\
\end{tabularx}

파란 상자 = 관제사 송신, 주황 상자 = 조종사 송신, 초록 막대 = 지시의 유효 구간 $\tau$. 시각은 2024-08-26 00:00:00Z부터의 초. NEXT는 실제로 \#0005의 마지막 지시(TRAFFIC\_INFO)에서 \#0022 MONITOR로 이어진다.}
\newpage
\section*{후보 D: 복창을 지시 단위로 연결 (구조 변경)}
READBACK\_OF를 송신끼리가 아니라 \textbf{조종사 지시 $\to$ 관제사 지시}로 잇는다. 지금은 어느 지시가 맞았는지가 송신 간선의 속성(\texttt{matched}, \texttt{conflicts})에 들어 있는데, 이를 간선 자체로 만든다. 예: \#0006의 TAXI\_VIA $\to$ \#0005의 TAXI\_VIA \{conflict: target\}. \textbf{코드 수정(M4, M5)이 필요하다.}
\begin{center}
\resizebox{0.95\textwidth}{!}{%
\begin{tikzpicture}[
  font=\small,
  dyn/.style={draw, thick, rounded corners=2pt, fill=comp, minimum height=0.8cm, minimum width=2.4cm, align=center},
  sta/.style={draw, thick, fill=res, minimum height=0.8cm, minimum width=2.2cm, align=center},
  e/.style={-{Stealth[length=2.2mm]}, thick},
  el/.style={font=\scriptsize\sffamily, fill=white, inner sep=1pt}
]
\node[dyn] (tx) at (0,0) {\textbf{Transmission}\\{\scriptsize $s, w$, speaker}};
\node[sta] (role) at (-4.4,0) {\textbf{Role}\\{\scriptsize ATCO / PILOT}};
\node[dyn] (fl) at (4.8,0) {\textbf{Flight}\\{\scriptsize $c$ (예: DAL1489)}};
\node[sta] (al) at (9.0,0) {\textbf{Airline}\\{\scriptsize icao, country}};
\node[dyn] (cm) at (4.8,2.2) {\textbf{CallsignMention}\\{\scriptsize $\hat c$ (원 표면형)}};
\node[dyn] (ins) at (0,-3.0) {\textbf{Instruction}\\{\scriptsize $\kappa,v,q,\tau$, speaker}};
\node[sta] (rw) at (4.8,-3.0) {\textbf{Runway / Taxiway}\\{\scriptsize / Ramp / Facility}};
\node[sta] (ap) at (9.0,-3.0) {\textbf{Airport}};
\draw[e] (tx) -- node[el,above]{SPOKEN\_BY} (role);
\draw[e] (tx) -- node[el,above]{ADDRESSES} (fl);
\draw[e] (fl) -- node[el,above]{OPERATED\_BY} (al);
\draw[e] (fl) -- node[el,right]{ALIAS} (cm);
\draw[e] (tx) -- node[el,right]{CONTAINS} (ins);
\draw[e] ([yshift=2mm]ins.east) -- node[el,above]{TARGET} ([yshift=2mm]rw.west);
\draw[e] ([yshift=-2mm]ins.east) -- node[el,below]{VIA} ([yshift=-2mm]rw.west);
\draw[e] (rw) -- node[el,above]{PART\_OF} (ap);
\draw[e, red!60!black] (ins.south west) to[out=-110,in=180,looseness=1.2] ($(ins.south)+(-0.2,-1.0)$) node[el,below,text=red!60!black,align=center]{READBACK\_OF \{matched/conflict\}\\{\scriptsize 조종사 지시 $\to$ 직전 관제사 지시: 같은 항공편·종류, $0<\Delta t\le 20$ s}} to[out=0,in=-70,looseness=1.2] (ins.south);
\draw[e, teal!60!black] (ins.west) to[out=200,in=270,looseness=1.3] ($(ins.west)+(-1.6,0.6)$) node[el,left,text=teal!60!black,align=right]{NEXT \{$\Delta t$\}\\{\scriptsize 같은 항공편의}\\{\scriptsize 다음 관제사 지시}\\{\scriptsize (시간순, $\Delta t\ge 0$)}} to[out=90,in=160,looseness=1.3] ([yshift=3mm]ins.west);
\end{tikzpicture}}
\end{center}

\end{document}
